\begin{afterword}
\summaryhead

The post continues the previous day's example of the lawyer at the soup kitchen. The author tells of a small good deed, telling a household that their car's trunk was open, and asks whether someone who works in the nonprofit sector can defend such a minute as the best use of time.

The defense is that the deed restores the author's willpower and keeps the author altruistic. That makes it a ``selfish good deed,'' and the author reports, with surprise, that calling it selfish does not spoil the effect, though this may differ between people.

The advice follows: buy warm fuzzies, status and utilons separately. A billionaire should buy fuzzies with a personal, anonymous gift to someone in need, status with a showy donation, and utilons by cold expected-value calculation, spending at least 20 times as much on utilons as on fuzzies. Trying to buy all three at once does none of them well. Ordinary donors should get their fuzzies cheaply, at a soup kitchen or by holding doors, and for utilons ``shut up and multiply.''

\responsehead

The rule is clear, and it is useful: it lets a donor keep the pleasures of giving without mistaking them for its effects. It also has a precedent. In 1997 Steven Landsburg divided giving by the same two motives, caring about the recipients and feeling good about giving, and concluded that a donor who cares about recipients should concentrate on the worthiest charity. James Andreoni's models of the ``warm glow'' (1989, 1990) gave the second motive its name in economics. What the post adds is the advice to buy the feeling on purpose, cheaply, as a small overhead on giving for utilons. Its name for the result, a ``selfish good deed,'' agrees with what the author had written in ``Circular Altruism'' in January 2008.

The rule rests on two supports. The first is that fuzzy acts are worth buying at all: they restore willpower and keep one altruistic. For this the post offers the author's own experience, flagged as a possible ``rationalization,'' a suspicion (``I suspect''), and ``Your mileage may vary.'' As advice drawn from one's own case, that is fair, and it is labelled as such. Later research, given in the notes as information, points in different directions and settles nothing.

The second support is that buying separately is far more efficient, because the best source of each good yields little of the others. The paragraph that argues this rests on two estimates marked ``probably'': that a \$10 million check to a breast-cancer charity gives far less glow than turning one life around in person, and that ``at least a thousand'' charities could do ``orders of magnitude'' more good with the money. No figures are given.

The calculation the post prescribes for utilons takes one side of a live question: how far to trust explicit expected-value estimates when evidence is thin. GiveWell's Holden Karnofsky later argued for the other side, preferring strong evidence of doing a lot of good to weak evidence of doing far more. The post's ratio of 20 to 1 is offered as the author's judgment, and its arithmetic holds.

In short: a clear rule for giving, close to one Landsburg stated in 1997, resting on the author's own experience of willpower, offered as such, and on a hedged premise that the best sources of warm feelings and of good outcomes lie far apart.
\end{afterword}
